% document formatting 
\documentclass[10pt]{article}
\usepackage[utf8]{inputenc}
\usepackage[left=1in,right=1in,top=1in,bottom=1in]{geometry}
\usepackage[T1]{fontenc}
\usepackage{xcolor}

% math symbols, etc.
\usepackage{amsmath, amsfonts, amssymb, amsthm}
\usepackage{dsfont}

% lists
\usepackage{enumerate}
\usepackage{tabularx}
\usepackage{multirow}
\usepackage{array}
\usepackage{blkarray}

% images
\usepackage{graphicx} % for images
\usepackage{tikz}
\usetikzlibrary{matrix}

% code blocks
\usepackage{minted, listings} 

% verbatim greek
\usepackage{alphabeta}
\DeclareMathOperator*{\argmin}{arg\,min}
\DeclareMathOperator*{\argmax}{arg\,max}
\newcommand{\dd}{\text{d}}

\graphicspath{{./assets/images/}}

\title{02-620 Week 12 \\ \large{Machine Learning for Scientists}}
 
\author{Aidan Jan}

\date{\today}

\begin{document}
\maketitle

\section*{Neural Networks}
\subsection*{Overview}
Model: Neural Network Architectures:
\begin{itemize}
	\item Feedforward neural networks (aka. Fully Connected NN)
	\item Convolutional neural networks
	\item Graph neural networks
	\item Variational autoencoder
	\item Recurrent neural networks
\end{itemize}
Learning
\begin{itemize}
	\item Stochastic gradient descent
	\item Back propagation
\end{itemize}

\subsection*{Logistic Regression as Basic Neuron Unit}
Logistic regression for representing $f \::\: X \rightarrow Y$
\begin{itemize}
	\item $X = [X_1, \cdots, X_n]$ is a vector of input features
	\item $Y$ is boolean
\end{itemize}
\[P(Y = 1 | X = [X_1, \cdots, X_n]) = \frac{1}{1 + \exp(-\w_0 - \sum_i w_i X_i)}\]
\begin{center} 
	\includegraphics*[width=0.8\textwidth]{W12_1.png} 
\end{center}
\begin{itemize}
	\item Input: $X = [X_1, \cdots, X_n]$
	\item Output: $o = \frac{1}{1 + \exp(-w_0 - \sum_i w_i X_i)} = \sigma(w_0 + \sum_i w_i X_i)$
\end{itemize}




\end{document}